% !TEX root = ../main.tex
\chapter{贝叶斯更新与信息}
\label{ch:bayes}

\noindent\emph{用到的前置工具：条件期望（上一章，作为投影理解）；概率空间、$\sigma$-代数与测度（第二十四章）；条件概率与独立性（第二十七章）；正态分布的密度与配方。}
\medskip

经济学里有一整类问题，主角不是``价格''或``产量'',而是\textbf{信念}。一个买家不知道二手车是好是坏，但卖家的报价透露了一点风声；一个雇主看不见求职者的能力，但学历是个信号；中央银行不确定经济处在什么状态，但每个月的数据都在挤牙膏式地泄露真相；一个计量经济学家手上有参数 $\theta$ 的先验看法，而样本一到，看法就该随之调整。这些情形的共同结构是：决策者对某个未知量持有一个\emph{概率性的看法}，然后观察到一条信息，需要把看法\emph{更新}成一个新的概率分布。贝叶斯定理就是做这件事的唯一一致的算法。

本章按使用者的顺序展开：先把贝叶斯定理的离散与连续两个版本讲清楚——它们其实是同一句话的两种写法；再引入共轭先验，说明为什么某些先验能让更新变成纯代数、并算两个最常用的例子（Beta--二项、正态--正态）；接着用 $\sigma$-代数把``信息''本身形式化，说明``信息更多''精确地等于``$\sigma$-代数更细'';然后转向博弈论，看信号如何更新信念；最后点出一个深刻而朴素的事实——\textbf{后验信念是一个鞅}，也就是``你没法预先知道自己将被说服往哪边''。这些机制是不完全信息博弈、信息经济学、宏观学习与贝叶斯计量共同的发动机。

\section{贝叶斯定理：把看法翻过来}
\label{sec:bayes-1}

\subsection{离散形式}

贝叶斯定理没有任何新内容，它只是条件概率定义的一个重排。对两个事件 $A,B$（设 $\Prob{B}>0$），条件概率定义为
\[
    \Prob{A\given B}=\frac{\Prob{A\cap B}}{\Prob{B}} .
\]
把交集 $\Prob{A\cap B}$ 用两种方式拆开——$\Prob{A\given B}\Prob{B}$ 与 $\Prob{B\given A}\Prob{A}$——令二者相等，立刻得到

\thm{贝叶斯定理（离散）}{
设 $\set{H_1,\dots,H_k}$ 是样本空间的一个划分（互斥且穷尽，各 $\Prob{H_i}>0$），$E$ 为一个 $\Prob{E}>0$ 的事件。则对每个 $i$，
\[
    \Prob{H_i\given E}
    =\frac{\Prob{E\given H_i}\,\Prob{H_i}}{\Prob{E}}
    =\frac{\Prob{E\given H_i}\,\Prob{H_i}}{\sum_{j=1}^{k}\Prob{E\given H_j}\,\Prob{H_j}} .
\]
}

这里 $H_i$ 是``假设''（未知世界的诸种可能状态），$E$ 是观察到的``证据''。分母用的是\emph{全概率公式} $\Prob{E}=\sum_j\Prob{E\given H_j}\Prob{H_j}$，它把 $E$ 的总概率沿所有假设展开。这个公式的全部作用，是把你\emph{容易写下的方向}（在每个假设下证据出现的概率 $\Prob{E\given H_i}$）翻转成你\emph{真正想要的方向}（看到证据后每个假设的概率 $\Prob{H_i\given E}$）。

\state{贝叶斯定理的四个名字}{
把上式按比例读，是理解一切贝叶斯推断的关键：
\[
    \underbrace{\Prob{H_i\given E}}_{\text{后验}}
    \;\propto\;
    \underbrace{\Prob{E\given H_i}}_{\text{似然}}\,\times\,
    \underbrace{\Prob{H_i}}_{\text{先验}} .
\]
\textbf{先验}是看到证据前的看法，\textbf{似然}是``这个假设有多能解释已见到的证据'',\textbf{后验}是更新后的看法。分母 $\Prob{E}$ 不依赖 $i$，只是把右边归一化成一个概率分布的常数——所以更新的\emph{方向与相对强度}全在``似然 $\times$ 先验''里，分母只管``凑成 $1$''。记住``后验 $\propto$ 似然 $\times$ 先验''，比记那个带分母的完整公式有用得多。
}

\ex[疾病检测：为什么阳性不等于患病]{
某病在人群中患病率 $\Prob{H}=0.01$。一种检测对病人有 $99\%$ 的灵敏度（$\Prob{E\given H}=0.99$），对健康人有 $5\%$ 的假阳性率（$\Prob{E\given H^c}=0.05$）。一个人检测呈阳性（事件 $E$），他真的患病的概率 $\Prob{H\given E}$ 是多少？
}
\sol{
直接套公式：
\[
\ba
    \Prob{H\given E}
    &=\frac{\Prob{E\given H}\Prob{H}}
          {\Prob{E\given H}\Prob{H}+\Prob{E\given H^c}\Prob{H^c}}\\[4pt]
    &=\frac{0.99\times 0.01}{0.99\times 0.01+0.05\times 0.99}
    =\frac{0.0099}{0.0099+0.0495}\approx 0.167 .
\ea
\]
阳性者中真正患病的不到 $17\%$。直觉上的``$99\%$ 准''之所以骗人，是因为它忽略了先验：健康人是病人的 $99$ 倍，哪怕每个健康人只有 $5\%$ 概率误报，误报的\emph{绝对数量}（$0.0495$）仍远超真阳性（$0.0099$）。这正是先验的力量——它是贝叶斯更新里``忽略不得''的那一项，也是直觉最常出错的地方。
}

\subsection{连续形式：后验 $\propto$ 似然 $\times$ 先验}

经济学里未知量常常是连续的参数 $\theta$（一枚硬币的正面概率、一个分布的均值、需求弹性……）。这时``划分''变成 $\theta$ 的取值连续区间，概率质量换成概率密度，求和换成积分，但那句口诀一字不改。

\thm{贝叶斯定理（连续 / 密度形式）}{
设未知参数 $\theta$ 有先验密度 $\pi(\theta)$，数据 $x$ 在给定 $\theta$ 下的条件密度（即\emph{似然}）为 $f(x\mid\theta)=L(\theta;x)$。则 $\theta$ 的后验密度为
\[
    \pi(\theta\mid x)
    =\frac{f(x\mid\theta)\,\pi(\theta)}{\displaystyle\int f(x\mid\theta')\,\pi(\theta')\,\d\theta'}
    \;\propto\; L(\theta;x)\,\pi(\theta) .
\]
分母 $m(x)=\int f(x\mid\theta')\pi(\theta')\d\theta'$ 称为\emph{边际似然}（证据），是关于 $\theta$ 的常数。
}

\rmkb[为什么实践中可以盯着比例号不放]{
分母 $m(x)$ 往往是个难算的积分，但它\emph{不依赖} $\theta$。许多时候我们根本不必算它：先把 $L(\theta;x)\pi(\theta)$ 写出来，认出它正比于某个已知分布的密度核，再补上那个分布该有的归一化常数即可。下一节的共轭先验把这套``认核''的手法用到极致——更新退化成对参数做加减法。$m(x)$ 本身也不是废料：在模型比较里它就是``这个模型整体上多能解释数据''的度量（贝叶斯因子的成分）。
}

多个观测如何累积？若数据 $x_1,\dots,x_n$ 在给定 $\theta$ 下\emph{条件独立同分布}，似然就是连乘 $L(\theta)=\prod_{i=1}^n f(x_i\mid\theta)$。一个等价而更能说明问题的看法是：贝叶斯更新可以\emph{串行}进行——今天的后验恰好充当明天的先验。

\state{贝叶斯更新是可迭代的，且与数据次序无关}{
把前 $n-1$ 个观测处理完得到的后验 $\pi(\theta\mid x_{1:n-1})$ 当作先验，再乘上第 $n$ 个观测的似然 $f(x_n\mid\theta)$，得到的后验
\[
    \pi(\theta\mid x_{1:n})\;\propto\; f(x_n\mid\theta)\,\pi(\theta\mid x_{1:n-1})
\]
与一次性用全部数据 $\propto\prod_{i}f(x_i\mid\theta)\,\pi(\theta)$ 完全相同。于是信念的更新是\emph{递归}的（这正是卡尔曼滤波与宏观学习模型的运作方式：来一个新数据点，就把信念往前推一步），而且——因为乘法可交换——\textbf{最终后验不依赖数据到达的先后顺序}，只依赖见到了哪些数据。
}

\section{共轭先验：让更新变成代数}
\label{sec:bayes-2}

一般而言后验是个新的、未必有名字的分布，分母那个积分也未必算得出。但若先验选得``聪明'',似然乘上先验后仍落在\emph{同一个分布族}里，更新就只是把先验的参数挪一挪。

\defn{共轭先验}{
对一个给定的似然族 $\set{f(x\mid\theta)}$，若先验 $\pi(\theta)$ 与后验 $\pi(\theta\mid x)$ 属于同一个参数化分布族，则称该先验族为这一似然的\emph{共轭先验}族。
}

共轭性不是天降的巧合，而是似然函数作为 $\theta$ 的函数其``形状''被先验复刻的结果：只要先验的密度核与似然的核是同一种代数形式，二者相乘后形式不变，更新自然退化为参数的加减。下面两例是全经济学用得最多的。

\subsection{Beta--二项：更新就是数数}

设未知量是一枚硬币（或``成功率''、``违约率''、``某政策被支持的比例'')的正面概率 $\theta\in[0,1]$。先验取 $\mathrm{Beta}(\alpha,\beta)$，其密度核为
\[
    \pi(\theta)\propto \theta^{\alpha-1}(1-\theta)^{\beta-1} .
\]
观测 $n$ 次独立投掷，得到 $s$ 次正面、$n-s$ 次反面。二项似然的核（把 $\binom{n}{s}$ 这种不含 $\theta$ 的常数丢进比例号）是
\[
    L(\theta)\propto \theta^{s}(1-\theta)^{n-s} .
\]
两者相乘：
\[
    \pi(\theta\mid s)\;\propto\;
    \theta^{s}(1-\theta)^{n-s}\cdot\theta^{\alpha-1}(1-\theta)^{\beta-1}
    =\theta^{(\alpha+s)-1}(1-\theta)^{(\beta+n-s)-1} .
\]
这正是 $\mathrm{Beta}(\alpha+s,\ \beta+n-s)$ 的核。于是

\prop{Beta--二项更新}{
若先验为 $\theta\dist\mathrm{Beta}(\alpha,\beta)$，观测到 $s$ 次成功、$f=n-s$ 次失败，则后验为
\[
    \theta\mid\text{数据}\;\dist\;\mathrm{Beta}(\alpha+s,\ \beta+f) .
\]
}

更新法则美到近乎透明：\textbf{见一次成功，把第一个参数加一；见一次失败，把第二个参数加一}。先验参数 $(\alpha,\beta)$ 因此可以解读为``伪计数''——仿佛你在看真数据之前，已经``虚拟地''见过 $\alpha-1$ 次成功、$\beta-1$ 次失败。后验均值
\[
    \E{\theta\mid\text{数据}}=\frac{\alpha+s}{\alpha+\beta+n}
\]
是先验均值 $\tfrac{\alpha}{\alpha+\beta}$ 与样本频率 $\tfrac{s}{n}$ 的加权平均；数据越多（$n$ 越大），权重越偏向数据，后验也越集中。图~\ref{fig:bayes-beta} 把这个``越学越窄、众数随证据移动''的过程画了出来。

\begin{figure}[h]
\centering
\begin{tikzpicture}[>=Stealth]
  % axes
  \draw[->] (-0.2,0) -- (9.6,0) node[right] {$\theta$};
  \draw[->] (0,-0.2) -- (0,4.4) node[above] {密度};
  \foreach \xv/\lab in {0/0,2.25/0.25,4.5/0.5,6.75/0.75,9/1}{
    \draw (\xv,0.06) -- (\xv,-0.06) node[below] {\small\lab};
  }
  % Beta(2,2) prior
  \draw[thick, gray!70!black] plot[smooth] coordinates {(0.000,0.000) (0.450,0.271) (0.900,0.513) (1.350,0.727) (1.800,0.912) (2.250,1.069) (2.700,1.197) (3.150,1.297) (3.600,1.368) (4.050,1.411) (4.500,1.425) (4.950,1.411) (5.400,1.368) (5.850,1.297) (6.300,1.197) (6.750,1.069) (7.200,0.912) (7.650,0.727) (8.100,0.513) (8.550,0.271) (9.000,0.000)};
  % Beta(5,3) after 3H1T
  \draw[thick, blue!55!black] plot[smooth] coordinates {(0.000,0.000) (0.900,0.008) (1.350,0.036) (1.800,0.102) (2.250,0.219) (2.700,0.396) (3.150,0.632) (3.600,0.919) (4.050,1.237) (4.500,1.559) (4.950,1.848) (5.400,2.068) (5.850,2.181) (6.000,2.189) (6.300,2.155) (6.750,1.973) (7.200,1.634) (7.650,1.172) (8.100,0.654) (8.550,0.203) (8.850,0.026) (9.000,0.000)};
  % Beta(14,6) after 12H4T
  \draw[thick, red!65!black] plot[smooth] coordinates {(0.000,0.000) (3.000,0.013) (3.450,0.053) (3.900,0.172) (4.350,0.447) (4.800,0.967) (5.250,1.759) (5.700,2.704) (6.000,3.270) (6.300,3.641) (6.450,3.715) (6.600,3.699) (6.900,3.382) (7.200,2.721) (7.500,1.859) (7.800,1.014) (8.100,0.393) (8.400,0.083) (8.700,0.004) (9.000,0.000)};
  % labels
  \node[gray!70!black, anchor=east] at (3.3,0.55) {先验 $\mathrm{Beta}(2,2)$};
  \node[blue!55!black, anchor=west] at (6.75,2.55) {$+\,3\text{ 正 }1\text{ 反}$};
  \draw[blue!55!black,->] (6.7,2.5) -- (6.1,2.25);
  \node[red!65!black, anchor=west] at (7.55,3.55) {$+\,12\text{ 正 }4\text{ 反}$};
  \draw[red!65!black,->] (7.5,3.55) -- (6.75,3.7);
\end{tikzpicture}
\caption{Beta--二项更新中后验密度的演化。先验 $\mathrm{Beta}(2,2)$ 平而宽；随观测增多，后验既向数据所示的频率（这里正面偏多）移动，又越来越窄——证据越多，信念越笃定。三条曲线的众数分别约在 $0.5,\,0.67,\,0.72$。}
\label{fig:bayes-beta}
\end{figure}

\rmkb[无信息先验与拉普拉斯接续律]{
取 $\alpha=\beta=1$ 即 $[0,1]$ 上的均匀分布（``无信息''先验）。此时观测 $s$ 次成功、$n-s$ 次失败后，后验均值 $\tfrac{s+1}{n+2}$ 正是历史上著名的``拉普拉斯接续律''。它给``到目前为止全是成功''（$s=n$）的情形一个不为 $1$ 的预测 $\tfrac{n+1}{n+2}$——朴素频率会鲁莽地断言``下次必成'',贝叶斯则因先验的存在而留有余地。
}

\subsection{正态--正态：精度相加、均值加权}

设未知量是一个连续参数 $\mu$（某资产的真实预期收益、某地区的真实平均工资、状态变量的真值），观测带有正态噪声。这是宏观学习与卡尔曼滤波的核心结构。

\asm{正态--正态模型}{
先验 $\mu\dist\cN(\mu_0,\tau_0^2)$。给定 $\mu$，观测一个数据点 $y\dist\cN(\mu,\sigma^2)$，其中观测方差 $\sigma^2$ 已知。
}

写出后验的核（把不含 $\mu$ 的因子统统并入比例号）：
\[
    \pi(\mu\mid y)\;\propto\;
    \exp\!\br{-\frac{(y-\mu)^2}{2\sigma^2}}
    \exp\!\br{-\frac{(\mu-\mu_0)^2}{2\tau_0^2}} .
\]
指数上是关于 $\mu$ 的二次式，所以后验必是正态；剩下的只是\emph{配方}找出它的均值与方差。把两个二次式中的 $\mu^2$ 与 $\mu$ 项归并：
\[
    -\frac12\br{\pr{\frac{1}{\sigma^2}+\frac{1}{\tau_0^2}}\mu^2
    -2\pr{\frac{y}{\sigma^2}+\frac{\mu_0}{\tau_0^2}}\mu}+\text{常数} .
\]
与 $\cN(\mu_1,\tau_1^2)$ 的核 $-\tfrac{1}{2\tau_1^2}(\mu-\mu_1)^2=-\tfrac{1}{2}\big(\tfrac{1}{\tau_1^2}\mu^2-\tfrac{2\mu_1}{\tau_1^2}\mu\big)+\text{常数}$ 逐项比对系数，立刻读出后验参数。引入\emph{精度}（方差的倒数）$\rho\equiv 1/\sigma^2$、$\rho_0\equiv 1/\tau_0^2$ 后，结论格外干净：

\prop{正态--正态更新}{
在上述假设下，后验为 $\mu\mid y\dist\cN(\mu_1,\tau_1^2)$，其中
\[
    \underbrace{\frac{1}{\tau_1^2}}_{\rho_1}
    =\underbrace{\frac{1}{\tau_0^2}}_{\rho_0}+\underbrace{\frac{1}{\sigma^2}}_{\rho},
    \qquad
    \mu_1=\frac{\rho_0\,\mu_0+\rho\,y}{\rho_0+\rho}
    =\frac{\tau_0^{-2}\mu_0+\sigma^{-2}y}{\tau_0^{-2}+\sigma^{-2}} .
\]
}

这两个式子值得刻进脑子。\textbf{精度相加}：后验精度等于先验精度加数据精度——信息是可叠加的``确定性''。\textbf{均值是按精度加权的平均}：后验均值落在先验均值 $\mu_0$ 与观测 $y$ 之间，谁的精度高就更靠近谁。先验越散（$\tau_0^2$ 大、$\rho_0$ 小），后验越被数据牵着走；观测噪声越大（$\sigma^2$ 大、$\rho$ 小），后验越守着先验不动。图~\ref{fig:bayes-normal} 画的正是这件事：后验夹在先验与似然之间、且比两者都更尖。

\begin{figure}[h]
\centering
\begin{tikzpicture}[>=Stealth, font=\small]
  \draw[->] (-0.2,0) -- (9.3,0) node[right] {$\mu$};
  \draw[->] (0,-0.2) -- (0,3.8) node[above] {密度};
  \draw (3.3,0.06) -- (3.3,-0.06) node[below] {\small$\mu_0$};
  \draw (5.5,0.06) -- (5.5,-0.06) node[below] {\small$y$};
  % prior N(0,1)  peak (3.3,1.835)
  \draw[thick, gray!70!black] plot[smooth] coordinates {(0.000,0.020) (0.440,0.062) (0.880,0.163) (1.320,0.363) (1.760,0.689) (2.200,1.113) (2.640,1.533) (3.080,1.799) (3.300,1.835) (3.520,1.799) (3.960,1.533) (4.400,1.113) (4.840,0.689) (5.280,0.363) (5.720,0.163) (6.160,0.062) (6.600,0.020) (7.040,0.006) (7.480,0.001) (8.000,0.000)};
  % likelihood N(2,0.5)  peak (5.5,2.595)
  \draw[thick, blue!55!black, densely dashed] plot[smooth] coordinates {(2.640,0.003) (3.080,0.021) (3.520,0.102) (3.960,0.366) (4.400,0.955) (4.840,1.811) (5.060,2.212) (5.280,2.494) (5.500,2.595) (5.720,2.494) (5.940,2.212) (6.160,1.811) (6.600,0.955) (7.040,0.366) (7.480,0.102) (7.920,0.021) (8.360,0.003) (8.800,0.000)};
  % posterior N(4/3,1/3)  peak (4.767,3.179)
  \draw[very thick, red!65!black] plot[smooth] coordinates {(2.567,0.008) (2.933,0.049) (3.300,0.221) (3.667,0.709) (4.033,1.632) (4.400,2.691) (4.583,3.049) (4.767,3.179) (4.950,3.049) (5.133,2.691) (5.500,1.632) (5.867,0.709) (6.233,0.221) (6.600,0.049) (6.967,0.008) (7.333,0.001) (7.700,0.000)};
  \draw[red!65!black, dotted] (4.767,0) -- (4.767,3.179);
  \fill[red!65!black] (4.767,0) circle (1.3pt);
  \node[red!65!black, anchor=north] at (4.767,-0.12) {\small$\mu_1$};
  \node[gray!70!black, anchor=east] at (2.5,1.4) {先验};
  \node[blue!55!black, anchor=west] at (6.0,2.1) {似然（数据）};
  \node[red!65!black, anchor=west] at (5.2,3.0) {后验};
\end{tikzpicture}
\caption{正态--正态更新。后验均值 $\mu_1$ 是先验均值 $\mu_0$ 与观测 $y$ 的精度加权平均，必落在二者之间（此处数据精度高于先验，故 $\mu_1$ 更靠近 $y$）；后验精度是二者之和，故后验比先验和似然都更尖。}
\label{fig:bayes-normal}
\end{figure}

\rmkb[这就是卡尔曼滤波的一步]{
把上式中的``一个观测''换成一批观测的均值 $\bar y$（其精度为 $n/\sigma^2$），再让 $\mu$ 随时间按某个状态方程演化，把``今天的后验''当``明天的先验''反复迭代，正态--正态更新就长成了\textbf{卡尔曼滤波}：每一步都是``预测—观测—按精度加权地修正''。宏观学习模型（agent 用过去数据估计未知参数、再据此决策）的递归核心，往往就是这个公式。
}

\section{信息的语言：$\sigma$-代数}
\label{sec:bayes-3}

到此我们一直在更新关于参数的分布。但``信息''本身——你\emph{知道}什么、能\emph{分辨}什么——也需要一个数学对象来承载，尤其在博弈论与动态模型里，``谁在何时知道什么''是问题的全部。这个对象就是 $\sigma$-代数（第二十四章已作为可测性的载体引入；这里给它一个``信息''的解读）。

直觉是这样的：你掌握的信息，等于你能回答的所有``是非问句''的集合。能不能分辨``状态落在集合 $A$ 里还是 $A$ 外''？若能，$A$ 就是你信息结构中的一个``可知事件''。所有这些可知事件构成一个对补集、可数并封闭的集类——一个 $\sigma$-代数。

\state{$\sigma$-代数就是``可分辨的事件''之集}{
把 $\sigma$-代数 $\cF\se 2^{\Omega}$ 读作``持有信息 $\cF$ 的人能够分辨的事件全体''。$A\in\cF$ 意味着此人能判断真实状态是否落在 $A$ 中。封闭性恰好对应推理的封闭性：若你能分辨 $A$，就能分辨它的反面 $A^c$（否定）；若你能分辨 $A$ 与 $B$，就能分辨 $A\cup B$（析取）。一个随机变量 $X$ 关于 $\cF$ \emph{可测}，正是说``知道 $\cF$ 的人就知道 $X$ 的取值''。
}

由此，``信息多少''有了精确的比较标准：

\defn{信息的精粗（$\sigma$-代数的细化）}{
设 $\cG,\cF$ 都是 $\Omega$ 上的 $\sigma$-代数。若 $\cG\se\cF$，则称 $\cF$ 比 $\cG$ \emph{更细}（$\cG$ 比 $\cF$ \emph{更粗}）。此时凡 $\cG$ 能分辨的事件 $\cF$ 都能分辨，故``$\cF$ 携带的信息不少于 $\cG$''。
}

最干净的图景是把信息看成对 $\Omega$ 的一个\emph{划分}：你知道真实状态落在划分的哪一块，但块内无法再细分。信息变多，等于把某些块进一步切碎——划分变细，对应的 $\sigma$-代数随之变细，能分辨的事件随之变多（图~\ref{fig:bayes-sigma}）。

\begin{figure}[h]
\centering
\begin{tikzpicture}[>=Stealth, font=\small]
  % ===== 左：粗 σ-代数（两块）=====
  \begin{scope}
    \draw[thick] (0,0) rectangle (4,3.2);
    \draw[thick] (2,0) -- (2,3.2);
    \node at (1,1.6) {$A_1$};
    \node at (3,1.6) {$A_2$};
    \node at (2,-0.55) {粗：$\mathcal{G}=\sigma\{A_1,A_2\}$};
    \node[gray!60!black] at (1,-1.1) {（只知道在左半还是右半）};
  \end{scope}
  % 箭头：细化
  \draw[->, thick, violet!60!black] (4.5,1.6) -- (6.0,1.6)
     node[midway, above] {细化};
  \node[violet!60!black] at (5.25,1.05) {更多信息};
  % ===== 右：细 σ-代数（四块，是左边的加细）=====
  \begin{scope}[xshift=6.5cm]
    \draw[thick] (0,0) rectangle (4,3.2);
    \draw[thick] (2,0) -- (2,3.2);
    \draw[thick] (0,1.6) -- (4,1.6);
    \node at (1,2.4) {$B_1$};
    \node at (3,2.4) {$B_2$};
    \node at (1,0.8) {$B_3$};
    \node at (3,0.8) {$B_4$};
    \node at (2,-0.55) {细：$\mathcal{F}=\sigma\{B_1,\dots,B_4\}$};
    \node[gray!60!black] at (2,-1.1) {（四块都能分辨）};
  \end{scope}
\end{tikzpicture}
\caption{信息越多 $\Rightarrow$ 划分越细 $\Rightarrow$ $\sigma$-代数越细 $\Rightarrow$ 可分辨的事件越多。右图保留了左图的分界（$\cG\se\cF$），只是把每块进一步切开。条件期望 $\E{X\given\cF}$ 就是``在每一块上取平均''，块越细，平均越逼近 $X$ 本身。}
\label{fig:bayes-sigma}
\end{figure}

这套语言把上一章的条件期望接了回来。$\E{X\given\cF}$ 是``在 $\cF$ 这个信息下对 $X$ 的最优预测''：在划分的每一块上把 $X$ 抹平成块内均值。$\cF$ 越细（信息越多），每块越小，块内均值越贴近 $X$ 的真值——预测越准。在两个极端上：$\cF=\set{\varnothing,\Omega}$（什么都不知道）时 $\E{X\given\cF}=\E{X}$ 退化成无条件均值；$\cF$ 细到 $X$ 关于它可测（完全知道 $X$）时 $\E{X\given\cF}=X$。``学习''在这套语言里，就是 $\sigma$-代数随时间逐步变细的过程。

\section{信号与信念更新：博弈论的视角}
\label{sec:bayes-4}

不完全信息博弈里，参与者对某个未知量（对手的``类型''、世界的``状态'')持有先验信念，再观察到一个与该未知量相关的\emph{信号}，据此更新信念后行动。贝叶斯定理正是这里``更新''二字的定义。

\defn{信号结构与后验信念}{
设状态 $\theta$ 取值于有限集 $\Theta$，先验为 $\pi(\theta)$。一个\emph{信号}（或信息结构）由一族条件分布 $\set{p(s\mid\theta)}_{\theta\in\Theta}$ 给出：在状态 $\theta$ 下收到信号 $s$ 的概率。观察到 $s$ 后的后验信念为
\[
    \mu(\theta\mid s)=\frac{p(s\mid\theta)\,\pi(\theta)}{\sum_{\theta'\in\Theta}p(s\mid\theta')\,\pi(\theta')} .
\]
}

这与离散贝叶斯定理一字不差，换了套博弈论的说法而已：状态 $\theta$ 是``假设'',信号 $s$ 是``证据'',$p(s\mid\theta)$ 是似然。信号的``信息量''可以用它在多大程度上改变信念来衡量——一个与 $\theta$ 完全独立的信号（$p(s\mid\theta)$ 不依赖 $\theta$）让后验等于先验，毫无信息；一个能完美区分状态的信号（不同 $\theta$ 给出不相交的 $s$ 支撑）让后验退化成一点，揭示全部真相。绝大多数现实信号介于两者之间。

\ex[一个二值信号如何更新信念]{
状态 $\theta\in\set{1,0}$（``高/低''），先验 $\pi(1)=\pi(0)=\tfrac12$。信号 $s\in\set{g,b}$ 满足 $p(g\mid 1)=0.8$、$p(g\mid 0)=0.2$（高状态更可能发出``好''信号）。收到 $s=g$ 后，关于 $\theta=1$ 的后验信念是多少？
}
\sol{
\[
    \mu(1\mid g)=\frac{p(g\mid 1)\pi(1)}{p(g\mid 1)\pi(1)+p(g\mid 0)\pi(0)}
    =\frac{0.8\times 0.5}{0.8\times 0.5+0.2\times 0.5}
    =\frac{0.4}{0.5}=0.8 .
\]
信念从 $0.5$ 升到 $0.8$。对称地，收到 $s=b$ 时 $\mu(1\mid b)=\tfrac{0.2\times0.5}{0.2\times0.5+0.8\times0.5}=0.2$，信念降到 $0.2$。信号把一个``五五开''的先验，分裂成``八二开''或``二八开''两种后验——信号越精确（$p(g\mid1)$ 与 $p(g\mid0)$ 差得越开），分裂得越厉害。
}

\state{这就是贝叶斯纳什均衡的``信念''部件}{
不完全信息博弈的均衡（贝叶斯纳什均衡、完美贝叶斯均衡）由两件相互咬合的东西组成：(i) 给定信念，每个类型的策略是最优的；(ii) 信念由先验加上``沿均衡路径观察到的行动/信号''按贝叶斯定理更新而来。本节的更新公式正是 (ii)。信号传递博弈（如教育作为能力信号）、筛选（如保险菜单甄别风险类型）、信息设计 / 贝叶斯劝说（设计 $p(s\mid\theta)$ 去影响接收者的后验、进而影响其行动），骨架全是``先验 + 信号 $\to$ 后验''这一步。
}

\section{后验信念是鞅：你说服不了自己往哪边走}
\label{sec:bayes-5}

本章最后一个事实，初看玄妙，证起来只是一行，却管着``学习''能与不能做到什么。问题是这样的：明天会有一个信号到来，把今天的信念 $\pi$ 更新成某个后验。\emph{在看到信号之前}，我对明天那个后验的期望是多少？答案是：还是 $\pi$。

\prop{后验信念的鞅性}{
设今天的信念为 $\pi(\theta)=\Prob{\theta}$，明天将观察到信号 $s$（其在状态 $\theta$ 下的分布为 $p(s\mid\theta)$）。则对每个状态 $\theta$，以\emph{今天}的信息取期望，明天的后验信念的期望等于今天的信念：
\[
    \E{\mu(\theta\given s)}
    =\sum_{s}\Prob{s}\,\mu(\theta\mid s)=\pi(\theta) .
\]
}
\pf{
按全概率公式，今天看明天，信号 $s$ 出现的（边际）概率是 $\Prob{s}=\sum_{\theta'}p(s\mid\theta')\pi(\theta')$，它正是贝叶斯公式的分母。于是
\[
    \sum_s \Prob{s}\,\mu(\theta\mid s)
    =\sum_s \Prob{s}\cdot\frac{p(s\mid\theta)\pi(\theta)}{\Prob{s}}
    =\sum_s p(s\mid\theta)\,\pi(\theta)
    =\pi(\theta)\sum_s p(s\mid\theta)
    =\pi(\theta) ,
\]
最后一步用了 $\sum_s p(s\mid\theta)=1$（信号分布对 $s$ 求和为一）。$\Prob{s}$ 在分子分母约掉，是这个``回到先验''的恒等式的全部机关。
}

把这个结论沿时间串起来：随着信号一个个到来，信念序列 $\set{\pi_t}$ 满足 $\E{\pi_{t+1}\mid\text{今天信息}}=\pi_t$——这正是\textbf{鞅}的定义（详见后文``鞅''一章）。它的含义掷地有声：

\state{信念无法被系统性地、可预测地移动}{
鞅性说的是：\emph{你不可能预先知道一个未来信号会把你的信念往哪个方向推}。若你今天就能断定``明天的证据多半会让我更信 $\theta=1$'',那你今天的信念本就该更高——否则就违反了 $\E{\pi_{t+1}}=\pi_t$。换句话说，理性信念的变动只能是``意外'',不能是``可预测的漂移''。这条原理直接派生出若干著名结论：金融里资产价格（作为对未来现金流的理性信念）近似为鞅，故``可预测的超额收益''难以存在（有效市场的形式表达）；信息设计里``发送者无法靠信号系统性抬高接收者对某状态的平均信念''（贝叶斯劝说的 Bayes-plausibility 约束，正是这条鞅性）。能被预测地移动的信念，等于桌上没捡的钱。
}

\rmkb[鞅性不是说``信念不动'']{
注意鞅性约束的是后验的\emph{期望}，不是后验本身。明天的信念几乎一定会变（信号一来，$\mu(\theta\mid s)$ 通常 $\neq\pi$）；鞅性只说这些变动\emph{平均起来}为零、方向无法预知。学习仍在发生、信念仍在收敛——只是每一步的``惊讶''事前不可预测。这与上一节图~\ref{fig:bayes-sigma} 的 $\sigma$-代数细化是同一枚硬币的两面：信息在增加（$\sigma$-代数在变细、后验在变窄），但增量是不可预测的。
}

下图把鞅性具象化：先验是一个点，信号把它分裂成两个后验，而两个后验按各自出现的概率加权平均，又回到先验那个点。

\begin{figure}[h]
\centering
\begin{tikzpicture}[>=Stealth, font=\small]
  \draw[->] (0,-0.3) -- (0,5.0) node[above] {信念 $\Prob{\theta=1\mid\cdot}$};
  \foreach \yv/\lab in {0/0,2/0.5,4/1}{
    \draw (-0.06,\yv) -- (0.06,\yv) node[left, xshift=-3pt] {\small\lab};
  }
  \draw[dotted, gray] (0,4) -- (8.6,4);
  \draw[dotted, gray] (0,0) -- (8.6,0);
  \node[anchor=north] at (1.5,-0.55) {收到信号前};
  \node[anchor=north] at (7.0,-0.55) {收到信号后};
  % prior node at (1.5, 2.0)
  \fill[violet!60!black] (1.5,2.0) circle (2.2pt);
  \node[violet!60!black, anchor=south] at (1.5,2.18) {先验 $\pi=0.5$};
  % posterior nodes
  \coordinate (G) at (7.0,3.2);
  \coordinate (B) at (7.0,0.8);
  \fill[blue!60!black] (G) circle (2.2pt);
  \fill[red!65!black] (B) circle (2.2pt);
  \draw[thick, blue!60!black] (1.5,2.0) -- (G) node[pos=0.62, above, sloped] {好信号};
  \draw[thick, red!65!black] (1.5,2.0) -- (B) node[pos=0.62, below, sloped] {坏信号};
  \node[blue!60!black, anchor=south east] at (3.0,2.50) {$q_g=0.5$};
  \node[red!65!black, anchor=north east] at (3.0,1.50) {$q_b=0.5$};
  \node[blue!60!black, anchor=west] at (7.2,3.2) {后验 $=0.8$};
  \node[red!65!black, anchor=west] at (7.2,0.8) {后验 $=0.2$};
  \draw[->, violet!55!black] (8.4,3.2) to[bend left=20] (8.4,0.8);
  \node[violet!55!black, anchor=west, align=left] at (8.55,2.0)
    {$0.5\!\cdot\!0.8+0.5\!\cdot\!0.2$\\[-1pt]$=0.5=\pi$};
\end{tikzpicture}
\caption{后验信念的鞅性。一个二值信号把先验 $\pi=0.5$ 分裂成 $0.8$（概率 $0.5$）与 $0.2$（概率 $0.5$）两个后验；两者按出现概率加权平均，又精确回到 $\pi=0.5$。信念会变，但事前的期望不变——这正是 $\E{\pi_{t+1}\mid\text{今}}=\pi_t$。}
\label{fig:bayes-martingale}
\end{figure}

\section{它通向何处}
\label{sec:bayes-6}

贝叶斯更新是少数几个``一处学会、处处遇见''的机制之一。把本章的工具接回经济学主干：

\begin{itemize}
    \item \textbf{不完全信息博弈}。贝叶斯纳什均衡与完美贝叶斯均衡把``信念按贝叶斯定理更新''写进了均衡定义本身；信号传递（教育、广告）、筛选（保险、合同菜单）都是它的应用。
    \item \textbf{信息经济学与信息设计}。从 Akerlof 的柠檬市场到贝叶斯劝说，核心都是``谁能观察到什么信号、信号如何重塑后验''；上一节的鞅性（Bayes-plausibility）是信息设计可行性约束的来源。
    \item \textbf{宏观学习}。当 agent 不知道模型参数、只能边观察边学习时，``信念的演化''成为模型的状态变量；正态--正态更新与卡尔曼滤波是其标准工具，且``信念是鞅''限制了学习能制造怎样的动态。
    \item \textbf{贝叶斯计量}。把参数 $\theta$ 当随机的、给它先验、用数据算后验，是与频率派并行的一整套推断范式；后验均值（即条件期望，承接上一章）是常用的点估计，后验分位数给出可信区间。共轭先验让少数模型有解析后验，其余则靠 MCMC 从后验抽样。
    \item \textbf{动态规划中的信念状态}。部分可观测的动态决策问题（POMDP、带学习的最优停时）里，``当前后验信念''本身就是 Bellman 方程的状态变量，更新法则即本章的递归式（见后文动态规划部分）。
\end{itemize}

一条由条件概率定义重排而来的式子，串起了博弈论、信息经济学、宏观与计量。这再次印证全书的主线：认清一件工具的内核，等于一次拿到许多扇门的钥匙——而贝叶斯定理，是其中开门最多的那一把之一。
